\documentclass[11pt,a4paper]{article}
\usepackage[T1]{fontenc}
\usepackage{lmodern,amsmath,amssymb,amsthm,mathtools,booktabs,microtype,xcolor}
\usepackage[margin=26mm]{geometry}
\usepackage[colorlinks=true,linkcolor=blue!45!black,citecolor=blue!45!black,urlcolor=blue!45!black]{hyperref}
\newtheorem{theorem}{Theorem}[section]\newtheorem{lemma}[theorem]{Lemma}
\newtheorem{proposition}[theorem]{Proposition}
\newcommand{\Q}{\mathbb Q}\newcommand{\Z}{\mathbb Z}\newcommand{\HH}{\mathcal H}
\DeclareMathOperator{\Tr}{Tr}\DeclareMathOperator{\Nm}{Nm}
\setlength{\emergencystretch}{2em}
\title{A family of defining quartics in 226 variables}
\author{}\date{10 November 2026}
\begin{document}\maketitle
\begin{abstract}
Under the specified modular-realization, elliptic, pole-realization and
prime-pattern inputs, we give an integer quartic family in 226 total
variables whose sufficiently large members define the integers in the
rationals. The proved contact-height exponent $7/3$ becomes $7/6$ on
two disjoint contact blocks. A common elliptic chart uses four field
coordinates, and all elliptic comparison ratios use the plain formal
parameter. Five factors represent sixteen local tests. Clearing each
factor equation by an inverse fourth power removes five separately
stored field fourth powers. Three reciprocal weights per detector
suffice, and the local comparison order is three. An inverse weight
also gives a contact circuit with three scalar outputs per detector.
The quartic is a sum of 200 primitive quadratic squares, has 2835
monomials, and has coefficient height $149^{4N+2}+1$ for $N\ge53$.
The sufficient defining parameter remains ineffective because of the
additive constant in the aggregate trace-height theorem. No additional
$abc$ conjecture is used.
\end{abstract}

\section{The family and its arithmetic scope}
All variable counts include the distinguished input $X$. For a fixed
integer $N\ge1$ let $A_N=149^{2N+1}$. The complete tables
\begin{center}\small
\nolinkurl{quadratic-family226-v21.json}\\
\nolinkurl{polynomial-family226-v21.json}
\end{center}
specify 200 primitive integer
quadratic polynomials and their sum of squares
\[
 P_N(X,Y_1,\ldots,Y_{225})=\sum_{j=1}^{200}Q_{N,j}^2.
\]
The parameter $A_N$ is a coefficient, not a variable.

\begin{theorem}\label{thm:main}
Assume the arithmetic premises stated below. There is a fixed integer
$N_0\ge53$ such that every $N\ge N_0$ satisfies
\[
 x\in\Z\quad\Longleftrightarrow\quad
 \exists\mathbf y\in\Q^{225}\ P_N(x,\mathbf y)=0
 \qquad(x\in\Q).
\]
For every $N\ge1$ the polynomial $P_N$ has degree exactly four,
content one, and 2835 nonzero coordinate monomials. For $N\ge53$,
\[
                       H(P_N)=149^{4N+2}+1.
\]
The construction specifies every member exactly. The proof does not
determine a sufficient numerical value of $N_0$.
\end{theorem}

Set
\[
 K=\Q(\rho),\quad \rho^2=2,\quad D=(3-\rho)^2=11-6\rho,
 \qquad E:y^2=x^3-Dx,
\]
\[
 P=2(1+\rho,2\rho)
 =\left(\frac{57-28\rho}{8},\frac{-660+417\rho}{32}\right).
\]
We retain the following arithmetic premises, with the precise domains
and hypotheses recorded in \cite{Dependencies,Source,Weights}:
\begin{enumerate}
\item $2E(K)=\Z P$ and $h(x([n]P))\ge n^2-6$ for every nonzero
integer $n$, together with the specified point charts and group domains.
\item The modular-realization input of
\cite[Definition~6.1]{NewHeight}, with its normalized coefficient
character. The multiplicity theorem of
\cite[Theorem~4.1]{Multiplicity} proves
\begin{equation}\label{eq:baseheight}
 h(u)\le H M_5(u)^{7/3},\qquad H=2^{221234}.
\end{equation}
Here $M_5(u)$ is the odd contact radical outside $76176$ for
$\infty,0,1,2401,-343/32$, and $u$ avoids the four finite contacts.
Explicitly, if $u=r/s$ is reduced with $s>0$, then $M_5(u)$ is
the product of primes $p>76176$ having a positive odd valuation on
at least one of
\[
           s,\quad r,\quad r-s,\quad r-2401s,\quad32r+343s.
\]
The estimate is a consequence of the stated modular input, rather
than an additional $abc$ assumption.
\item The pole obstruction outside $2,3,17$, and its simultaneous
realization for positive inputs less than $17$ with square reduced
denominator and the stated local squareclasses and valuation bounds
at those three primes. Its normalized quadratic equations are retained.
\item The simultaneous prime-pattern, conjugate reduction-order,
finite congruence, and real approximation assertions used by the source.
The weighted selection and fixed denominator dilation are justified
from the cited analytic estimates in \cite{Selection}.
\end{enumerate}
The aggregate trace-height theorem used below is the proved consequence
of the family height theorem described in \cite{Aggregate}; its additive
constant is ineffective. The contact amplification below is an elementary consequence of~\eqref{eq:baseheight},
not an additional arithmetic assumption. Exact coefficient verification
does not prove the retained arithmetic premises.

\section{Two contact blocks and the comparison scales}
Put
\[
 \Theta=\{0,-1/2,-1,1/2399,-32/407\},\qquad
 \Theta_2=\Theta\cup(2-\Theta),\qquad
 F_2(t)=\prod_{\theta\in\Theta_2}(t-\theta).
\]
The ten contacts are distinct, omit one, and are symmetric about one.
Let $\Sigma$ consist of the primes at most $76176$ and the prime
$2029147$. Outside this set, all contact differences and all numbers
$1-\theta$ are units. Define $M_2^\circ(t)$ as the product of primes
$p\notin\Sigma$ for which $v_p(t-\theta)$ is positive and odd for
some $\theta\in\Theta_2$. Exact factorizations in the independent
certificate verify this exceptional set. In particular
\[
 F_2(1)=-\frac{82429855236}{7878870169},\qquad M_2^\circ(1)=1.
\]

Set $E_2=2029147$ and $H_2=2^{221284}$. Then
\begin{equation}\label{eq:contactheight}
                   h(t)\le H_2M_2^\circ(t)^{7/6}
                   \qquad(t\notin\Theta_2).
\end{equation}
Indeed, apply \eqref{eq:baseheight} to $2+1/t$ and $2+1/(2-t)$.
The inverse projective matrices have maximum absolute row sums three
and seven, respectively. If $M_+,M_-$ are the two block radicals
with $\Sigma$ omitted, the original radicals are at most $E_2M_+$
and $E_2M_-$. Outside $\Sigma$ these radicals are coprime, with
$M_+M_-=M_2^\circ(t)$. Multiplying the two inequalities
\[
             h(t)-\log7\le HE_2^{7/3}M_\pm^{7/3}
\]
when the left side is positive gives the exponent $7/6$.
If it is nonpositive, the asserted bound is immediate. Finally
$E_2<2^{21}$ and $\log7<2$ give
\[
                 \log7+HE_2^{7/3}<2+2^{221283}<H_2.
\]

The retained root points are
\[
 Q_1=[n]P,\quad Q_2=[a]P,\quad Q_3=[t]P,\quad Q_4=[b]P,
 \qquad B=\max(|n|,|a|,|t|,|b|),
\]
with nonzero integer labels. Set
\[
 T=\frac14\Tr_{K/\Q}\left(\sum_{i=1}^4x(Q_i)\right).
\]
It is positive, and the aggregate trace theorem gives fixed constants
$M_1>0$, $M_2\ge1$ such that
\begin{equation}\label{eq:aggregate}
                         B^2\le M_1h(T)+M_2.
\end{equation}
Use the two positive rational scales $a_*=1,T$. No square condition
on a scale is needed. Replacing the half-trace in the aggregate
height theorem by its half changes its logarithmic height by at
most $\log2$; this is absorbed into $M_2$.
For each scale put $t_*=a_*\lambda$ and $\delta=t_*-1$.
On the nonzero domain use the balanced contact and the local weight
\begin{equation}\label{eq:detector}
 R=\frac{F_2(t_*)}{\delta^{10}},\qquad
 q=\frac{\delta^{10}}{F_2(t_*)}-1=R^{-1}-1,
 \qquad \gamma=\frac{408}{409}.
\end{equation}
Let $e_j$ be the $j$th elementary symmetric sum of the five positive
rational numbers $(1-\theta)^2$, $\theta\in\Theta$. Thus
\[
 P_5(r)=\prod_{\theta\in\Theta}(r-(1-\theta)^2)
       =r^5-e_1r^4+e_2r^3-e_3r^2+e_4r-e_5,
 \qquad e_5\ne0.
\]
Use the inverse weight $s_1$ required by the local factors, with
$qs_1=1$. The contact equations retain exactly three scalar outputs
$r,r_2,v$:
\begin{equation}\label{eq:contactcircuit}
\begin{split}
 r&=\delta^2,\qquad r_2=r^2,\\
 v&=rr_2-e_1(1+s_1)r_2+e_2(1+s_1)r,\\
 r_2v&=(1+s_1)(e_3r_2-e_4r+e_5).
\end{split}
\end{equation}
These are four quadratic equations; the inverse weight is shared with
the factor equations below.

\begin{lemma}\label{lem:contactcircuit}
The equations $qs_1=1$ and~\eqref{eq:contactcircuit} imply
\begin{equation}\label{eq:clearing}
                       (1+q)P_5(r)=r^5,
\end{equation}
with $P_5(r)\ne0$. Conversely, any rational $r=\delta^2,q,s_1$
satisfying $qs_1=1$ and~\eqref{eq:clearing} has a unique extension
to $r_2,v$ satisfying~\eqref{eq:contactcircuit}. The case
$r=0$ is exactly $q=s_1=-1$, $r_2=v=0$.
\end{lemma}
\begin{proof}
Put $N(r)=r^5-P_5(r)$. Substituting the equations for $r_2,v$
into the last equation of~\eqref{eq:contactcircuit} gives
\[
                       r^5-(1+s_1)N(r)=0.
\]
Multiplication by $q$, using $qs_1=1$, gives $qP_5(r)=N(r)$,
which is~\eqref{eq:clearing}. Reversing these steps proves the
converse, and the displayed equations determine $r_2,v$ uniquely.
If $P_5(r)=0$, then~\eqref{eq:clearing} forces $r=0$,
contrary to $P_5(0)=-e_5\ne0$. At $r=0$, the last equation gives
$1+s_1=0$, and $qs_1=1$ gives $q=-1$. The other two outputs vanish.
\end{proof}

Since $P_5((t_*-1)^2)=F_2(t_*)$, the lemma excludes every contact
value. When $\delta\ne0$, it gives precisely $R^{-1}=1+q$.
The remaining stratum $\delta=0$ has $q=-1$ and
$M_2^\circ(t_*)=1$; no value of $R$ is used on that stratum.
Every witness constructed in completeness has a selected contact prime
with $v_p(q)=-1$, and therefore satisfies $q\ne0$.

If both radicals divide a nonzero integer $J$ to order $k$, then
$h(T)\le h(\lambda)+h(T\lambda)$ and
\eqref{eq:contactheight}--\eqref{eq:aggregate} imply
\begin{equation}\label{eq:heighttransfer}
 B^2\le C|J|^{7/(6k)}+M_2,\qquad C=2M_1H_2.
\end{equation}
The two radicals need not be coprime for this implication.

\section{A quadratic chart for the common reference point}

Write $a_P=x_P$, $b_P=y_P$, and set
\[
 A_P=a_P^2+D,\qquad B_P=a_P^2-D,\qquad
 C_P=D^2+6a_P^2D-3a_P^4.
\]
The constants $a_P,b_P,D,A_P,C_P$ are nonzero. The first four
nonvanishing assertions follow from the displayed coordinates and
$D=(3-\rho)^2>0$ at both real embeddings; exact arithmetic gives
\[
 C_P=-\frac{94897379}{4096}+\frac{4195821}{256}\rho,
 \qquad \Nm(C_P)=-\frac{8203357319351}{16777216}\ne0.
\]
Retain four field coordinates $S_H,T_H,f,F_0$. Put
\begin{align*}
 L_f&=2b_PD(a_PT_H-b_P),\\
 Q_f&=a_PA_P S_H T_H+b_PA_P S_H-b_PB_PT_H+4a_P^2D,\\
 L_g&=2b_PD(a_P S_H+b_P),\\
 Q_g&=-a_PA_P S_H T_H-b_PB_P S_H+b_PA_PT_H-4a_P^2D.
\end{align*}
Impose the three quadratic equations
\begin{equation}\label{eq:Hchart}
 \begin{split}
 L_ff&=Q_f,\\
 L_gF_0&=Q_g,\\
 D^2fF_0+a_PD(f+F_0)+a_PS_HT_H+b_P(S_H-T_H)+D&=0.
 \end{split}
\end{equation}

\begin{lemma}\label{lem:commonchart}
These equations describe exactly the points $H\in E(K)$ for which
$H$ and $H-P$ are finite with nonzero abscissas, through
\[
 S_H=\frac{y(H)}{x(H)},\qquad T_H=\frac{y(H-P)}{x(H-P)},\qquad
 f=\frac1{x(H)},\qquad F_0=\frac1{x(H-P)}.
\]
In particular $fF_0\ne0$ at every solution. The statement includes
all finite exceptional cases of the displayed equations.
\end{lemma}
\begin{proof}
For the algebraic calculation abbreviate $a=a_P$, $b=b_P$, $S=S_H$, $T=T_H$
and $g=F_0$. Thus $b^2=a(a^2-D)$. Define
\[
 \Phi=aS^2T^2-(a^2-D)(S^2+T^2)+2(a^2+D)ST+4aD
\]
and the polynomial expressions
\begin{align*}
 2bDf_0&=-(a^2+D)S+(a^2-D)T-bS^2-aS^2T,\\
 2bDg_0&=(a^2+D)T-(a^2-D)S-bT^2+aST^2.
\end{align*}
Direct expansion gives
\begin{align}
 Df_0^2+S^2f_0-1
   &=\frac{a^2S^2-b^2}{4ab^2D}\Phi,\label{eq:chartcurve}\\
 Dg_0^2+T^2g_0-1
   &=\frac{a^2T^2-b^2}{4ab^2D}\Phi,\nonumber\\
 L_ff_0-Q_f&=-a\Phi,\qquad L_gg_0-Q_g=a\Phi.\nonumber
\end{align}
Let $r_f=L_ff-Q_f$, $r_g=L_gg-Q_g$, and
$R=D^2fg+aD(f+g)+aST+b(S-T)+D$. The further exact identity
\begin{equation}\label{eq:chartmaster}
\begin{split}
 L_fL_gR={}&D^2r_fr_g+(aDL_g+D^2Q_g)r_f\\
          &+(aDL_f+D^2Q_f)r_g-aD^2C_P\Phi
\end{split}
\end{equation}
shows that every solution has $\Phi=0$, including the cases in which
one leading coefficient vanishes. Also
\[
 (2bD)^2R(f_0,g_0)=
 D^2\bigl((a^2-D)-aST-b(S-T)\bigr)\Phi.
\]
If $L_fL_g\ne0$, these identities give $f=f_0,g=g_0$.
If $L_f=0$, then $T=b/a$ and $\Phi=0$ forces
\[
 S=\frac{a^4-6Da^2+D^2}{2b(a^2+D)},\qquad
 L_g=-\frac{aDC_P}{a^2+D}\ne0.
\]
Thus $g=g_0=1/a$. The coefficient of $f$ in $R$ is
$D(a^2+D)/a\ne0$, so $f=f_0$. The case $L_g=0$ is symmetric.
In particular both leading coefficients cannot vanish simultaneously;
at $S=-b/a,T=b/a$ one has $Q_f=C_P\ne0$.

Equation~\eqref{eq:chartcurve} now proves $fg\ne0$ and gives the
points $H=(f^{-1},S/f)$ and $H'=(g^{-1},T/g)$ on $E$.
Two more identities prove the translation:
\[
 a(f_0T+g_0S)+b(g_0-f_0)=S+T,
\]
\[
 f_0g_0(aS+b)^2-a(1-af_0)^2
 =-\frac{(aS+b)^2\bigl(aST+b(S-T)+D\bigr)}{4b^2D^2}\Phi.
\]
If $x(H)\ne a$, the second identity gives the abscissa of
$H-P$ from the chord slope $(x(H)S+b)/(x(H)-a)$.
The first states that $H,-H',-P$ are collinear, and therefore gives
the required ordinate. If $x(H)=a$, then $S=\pm b/a$.
At $S=b/a$ the formula gives $f=-a/D\ne1/a$, which is impossible.
At $S=-b/a$ it gives $H=-P$ and
\[
 T=-\frac{a^4-6Da^2+D^2}{2b(a^2+D)},\qquad
 g=\frac{4b^2}{(a^2+D)^2},
\]
which are the coordinates of $H'=-2P$ in the stated form.
Thus $H'=H-P$ in all cases.

Conversely, the chord-and-tangent formulas for any allowed $H$ give
$\Phi=0$ and the expressions $f_0,g_0$ above; the same explicit
doubling computation covers $H=-P$. The displayed identities then
prove all three equations. This proves both directions.
\end{proof}

Write $P=2R_0$ and $H=[m]R_0+T_0$, with $T_0\in E[2](K)$, and put $h=m/2$.
This notation is an elliptic index; it does not define multiplication
of a point by a nonintegral rational number.


\section{The scalar certificate and the point domains}\label{sec:scalar}
Put $\beta=51X^2+2504$, let $S=Q_2=[a]P$, and write
$C_S=S-P=[a-1]P$. The scalar normalizer is
\begin{equation}\label{eq:scalarnormalizer}
 \begin{split}
 \ell&=149,\qquad A_N=\ell^{2N+1},\\
 r_s&=\beta(\beta-1)(\beta-2401)(32\beta+343)(2\beta-1),\\
 R_s&=\frac{A_Nr_s}{1+\ell^2r_s}.
 \end{split}
\end{equation}
All five factors of $r_s$ are positive, since $\beta\ge2504$, so
$R_s>0$. The retained square-abscissa chart for $S$ is
\[
 x_S=h_S^2,\qquad j_S^2=x_S^2-D,
 \qquad S=(x_S,h_Sj_S).
\]
It has $h_Sx_Sj_S\ne0$. The fixed-subtraction coordinate $g_S\in K$
satisfies
\[
 8(x_Px_S+D)g_S+(39-36\rho)h_S+(8-14\rho)j_S=0.
\]
The rational component of $g_S$ has a retained inverse witness.
At $S=P$ the coordinate $g_S$ is zero, so this guard gives $a\ne1$.
All nonzero points of $\Z P$ are nontorsion. Thus $C_S$ is a nonzero
nontorsion point, and the usual fixed-subtraction formulas have no
remaining zero denominator. Explicitly, define
\[
 \begin{split}
 N_S&=(-660+417\rho)h_S+(-228+112\rho)j_S,\\
 D_S&=(-704+384\rho)h_S+(660-417\rho)j_S
                         +(6888-5040\rho)g_S.
 \end{split}
\]
Then $D_S\ne0$ and $z(C_S)=N_S/D_S$, where $z=-x/y$.
These identities are obtained from the chord-and-tangent formulas;
they are the same fixed-subtraction identities independently checked
in the exact function field by the supplied certificate.

Use the common chart of Lemma~\ref{lem:commonchart}, so
$H=(f^{-1},S_H/f)$ and $H-P=(F_0^{-1},T_H/F_0)$.
Before any logarithm is used, the scalar ratio
coordinate $Z=Z_0+\rho Z_1$ is determined by
\[
                         ZD_S=-N_SS_H.
\]
When $H$ is nontorsion this is $Z=z(C_S)/z(H)$.
Set $F=\Nm_{K/\Q}(f)\ne0$. The reduced scalar norm equations are
\begin{equation}\label{eq:scalarnorm}
 A=\frac{Z_0^2+R_sZ_1^2}{R_s^4},\qquad
 W_s=\frac{A^2+R_sF^2}{R_s^4},\qquad
 b_s=R_s(R_s+9792)W_s^2.
\end{equation}
The actual quadratic source stores $R_s^2,R_s^4,Z_1^2,F^2,A,W_s$
and a scaled square of $W_s$; thus these equations are quadratic after
clearing. The normalized scalar pole has exactly the ordinary pole
input $b_s$ in~\eqref{eq:scalarnorm}. To verify its scaling literally,
let $M$ be the coefficient of the stored $W_s$-square coordinate in
source row491, let $K$ be the coefficient of $e\bar b$ in row499,
and let $L$ be the coefficient of $\bar b p^2$ in row494. The exact
coefficient identities are
\[
                  K=82479248M,\qquad 3L=17K.
\]
The rational substitutions
\[
 b_s=-K\bar b,\quad e'=-e/K,\quad p'=p/K,\quad
 u'=u/K,\quad v'=-v,\quad w'=w/K
\]
convert those rows to the retained normalized pole equations. The
independent checker verifies these identities. Since $R_s>0$ and
$F\ne0$, one has $W_s>0$ and $b_s>0$ at every rational source solution.

\begin{lemma}\label{lem:scalarextract}
Let $p>76176$ or $p=149$, and suppose $\nu=v_p(R_s)>0$ is odd. Then
\[
 v_p(F)>5\nu/4,\qquad
 v_p(Z_0)>23\nu/8,\qquad v_p(Z_1)>19\nu/8.
\]
Once the point interpretations hold, these imply
\[
                v_p(a-1)>19\nu/8.
\]
\end{lemma}
\begin{proof}
The number $9792=2^6\cdot3^2\cdot17$ is a unit at $p$.
The valuation of the nonzero pole input is
$\nu+2v_p(W_s)$, so the pole obstruction makes it positive and
$v_p(W_s)>-\nu/2$.
The two nonzero terms of a norm $u^2+R_sv^2$ have opposite
valuation parity. A zero term can be omitted. Hence the second
norm in~\eqref{eq:scalarnorm} gives
\[
 v_p(A)>7\nu/4,\qquad v_p(F)>5\nu/4.
\]
The first norm gives the two asserted bounds on $Z_0,Z_1$.

At some place above $p$ the value $f(H)$ has positive valuation,
because the valuation of its norm is positive. Thus $H$ is in the
formal group there. The displayed bounds give $v(Z)>19\nu/8>0$,
so $v(z(C_S))>0$. On $y^2=x^3-Dx$, at a good odd prime a point
with positive $z$-valuation is either formal or reduces to $(0,0)$.
Indeed $z^2=x/(x^2-D)$, and $D$ is a unit. Translation by $(0,0)$
sends $z$ to $-z$ and exchanges these two cases. Choose
$C'=C_S+\epsilon(0,0)$ to be formal. The integral identity
\[
                       [2m]C'=[4(a-1)]H
\]
kills both two-torsion terms. The formal logarithm is an isometry,
so its ratio has valuation $v_p((a-1)/h)=v(Z)$. Since $h=m/2$ is
$p$-integral, $v_p(a-1)\ge v(Z)>19\nu/8$.
\end{proof}

The integer $\ell=149$ is prime. For the five contact roots
$0,1,2401,-343/32,1/2$, the residues of $(\theta-2504)/51$ are
respectively $59,97,109,108,78$ modulo $149$. Their $74$th powers
are all $148$, so they are nonsquares. Trial division through $12$
certifies primality. The preceding lemma applies at $149$: the curve
has good reduction there, $9792$ is a unit, and the pole obstruction
excludes only $2,3,17$. The generic contact cutoff is unnecessary
at this separately verified prime. Thus $r_s$ is an $\ell$-adic unit when $X$ is integral
at $\ell$. If $v_\ell(X)<0$, then $v_\ell(r_s)=10v_\ell(X)$.
It follows that
\begin{equation}\label{eq:forcedvaluation}
 v_\ell(R_s)=
 \begin{cases}2N+1,&v_\ell(X)\ge0,\\2N-1,&v_\ell(X)<0.
 \end{cases}
\end{equation}
This is always positive and odd. These finite facts are included
in the independent exact certificate.

Lemma~\ref{lem:scalarextract}, using only its norm inequalities at
this stage, gives $v_\ell(F)>0$. Every finite allowed point in
\[
                E[2](K)\ \cup\ (P+E[2](K))\ \cup\ (-P+E[2](K))
\]
has reciprocal-abscissa norm an $\ell$-adic unit. To check this,
use the two-torsion translations and
\[
 \Nm(x_P)=41^2/64,\quad
 \Nm(x_P-(3-\rho))=17^2/64,\quad
 \Nm(x_P+(3-\rho))=63^2/64.
\]
The remaining torsion abscissas have norm seven. Thus none of
these points can equal $H$. In particular $h\notin\{-1,0,1\}$,
$H$ and $H-P$ are nontorsion, and all required denominators and
formal-logarithm interpretations are now justified. This exclusion
precedes every use of a comparison logarithm.

Let $c=|a-1|>0$. Outside the fixed cutoff, every prime of
$M_5(\beta)$ is a positive odd contact of $r_s$ and $R_s$.
The contact at infinity contributes none, since a negative valuation
of $\beta=51X^2+2504$ is even there. Lemma~\ref{lem:scalarextract}
therefore implies $M_5(\beta)^3\mid c$. The prime $\ell=149$ is below the defining cutoff of
$M_5$, so it is automatically absent from $M_5(\beta)$.
The forced-prime estimate is consequently independent of these radical
factors. We obtain
\begin{equation}\label{eq:scalardivisibility}
 \ell^{L_N}M_5(\beta)^3\mid c,\qquad
 L_N=\left\lfloor\frac{19(2N-1)}8\right\rfloor+1.
\end{equation}
Together with \eqref{eq:baseheight}, this gives
\begin{equation}\label{eq:scalar}
 h(\beta)\le H\left(\frac{c}{\ell^{L_N}}\right)^{7/9}
             \le H\ell^{-L_N}c,
 \qquad c\ge\ell^{L_N}.
\end{equation}
The last inequality uses $c/\ell^{L_N}\ge1$. These scalar estimates
are independent of the comparison logarithms.

\section{Plain parameter ratios and the local tests}
For $f=1/x$ and $z=-x/y$, extend their regular values at the
identity by zero. The retained point relations are
\[
 U_0=H,\quad V_0=H-P,\quad U_i=[u_i]P,\quad V_i=U_i-Q_i
 \qquad(1\le i\le4).
\]
The charts permit $U_i=O$ or $V_i=O$ for $i>0$.
For each numerator pair retain its quadratic pair-chart coordinates
$\mathcal V_i,\mathcal Q_i,\mathcal W_i,\mathcal I_i,\mathcal K_i$
as specified in~\cite{Even}. Their exact relations give
\[
 z(U_i)=-\frac{\mathcal V_i}{\mathcal I_i},\qquad
 z(V_i)=-\frac{\mathcal Q_i}{\mathcal K_i},\qquad
 \mathcal I_i\mathcal K_i\ne0,
\]
including the identity branches with zero numerator.
The common-point chart stores both negative reciprocal parameters
$S_H=-1/z(H)$ and $T_H=-1/z(H-P)$. Thus the quadratic equations
\[
 L_i\mathcal I_i=\mathcal V_iS_H,\qquad
 L_{i,0}'\mathcal K_i=\mathcal Q_iT_H,\qquad
 L_i'=\sigma L_{i,0}'
\]
mean precisely
\[
 L_i=\frac{z(U_i)}{z(H)},\qquad
 L_i'=\sigma\!\left(\frac{z(V_i)}{z(H-P)}\right).
\]
No numerator-parameter coordinate or logarithm correction is required.
Put
\[
 D_H=f(H),\quad A_i=f(U_i),\quad B_i=\sigma f(V_i),\quad
 C_i=L_i-L_i',\quad I=L_1-\beta,
\]
\[
 G=L_3^2-L_2(L_2-2)L_3L_4+(L_2-2)^2L_4^2-1.
\]
The required orders are two for $D_H$, three for $C_2,C_3,C_4,G$,
one for $C_1$, two for $I$, and one for the nine $A_i,B_i$.
The final input comparison uses order one.

\section{Five factors and extraction at one place}
At a contact prime $p\notin\Sigma$, $\delta$ is a unit and
$\nu=v_p(R)=v_p(F_2(t_*))$ is positive and odd, and $v_p(q)=-\nu$. This follows
from the unique contact factor with positive odd valuation and
does not require any condition on the valuation of the scale. All valuations
of $K$ are unramified and normalized by $v(p)=1$.
For a point on $E$, every finite value of $v(f)$ is even: if
$v(x)<0$ then $2v(y)=3v(x)$, while if $v(x)>0$ then
$2v(y)=v(x)$; the remaining case is valuation zero.
This also holds for conjugated reciprocal abscissas. The common-point lemma identifies $F_0=f(H-P)$ on every solution,
including the exceptional leading-coefficient cases. Its even valuation
is therefore available before the exclusion of torsion classes and
before any formal-parameter quotient is used.

Use
\begin{align*}
 T_1&=1+q^5C_2^2+q^2D_H^2+q^4B_1^4,\\
 T_2&=1+q^5G^2+q^2A_3^2+q^4B_0^4,\\
 T_3&=1+q^5I^4+q^3A_2^2+q^2A_4^2+q^4B_3^4,\\
 T_4&=1+q^5C_3^2+q^2C_1^4+q^4B_4^4,\\
 T_5&=1+q^5C_4^2+q^2A_1^2+q^4B_2^4,
 \qquad W=T_1T_2T_3T_4T_5.
\end{align*}

\begin{lemma}\label{lem:residues}
At a contact prime, every factor is nonzero and has valuation at most
zero. A factor of valuation zero is congruent to one in the residue field.
\end{lemma}
\begin{proof}
First suppose every displayed test in a factor is nonzero. In each
factor other than $T_3$, the three nonconstant valuations belong,
respectively, to the odd integers, to $2+4\Z$, and to $4+8\Z$.
For example $T_1$ has valuations
\[
 2v(C_2)-5\nu,\qquad 2v(D_H)-2\nu,\qquad4v(B_1)-4\nu.
\]
The second has exact factor two because $v(D_H)$ is even and $\nu$ is
odd. The third has exact factor four because $v(B_1)$ is even.
The term $q^2C_1^4$ in $T_4$ has valuation
$2(2v(C_1)-\nu)$, again twice an odd integer.

In $T_3$ the four nonconstant valuations are
\[
 4v(I)-5\nu,\quad 2v(A_2)-3\nu,\quad
 2v(A_4)-2\nu,\quad4v(B_3)-4\nu.
\]
The first three have residues $-\nu,\nu,2$ modulo four. These are
pairwise distinct and nonzero because $\nu$ is odd and the $A_i$
have even valuations. The fourth has residue four modulo eight;
it therefore differs from each of the first three and from zero.
Thus, in every factor, all finite term valuations, including zero
for the constant term, are distinct. Zero tests simply remove terms
from this argument. The minimum is attained exactly once. This proves
nonvanishing and the asserted bound. If the minimum is zero, every
nonconstant term has positive valuation, so the residue is one.
\end{proof}

\begin{proposition}\label{prop:extract}
At each contact prime, the retained pole system and the cleared trace
equations imply $\nu>0$ and give all sixteen required test orders at
one place above that prime.
\end{proposition}
\begin{proof}
Write $W=w_0+w_1\rho$. The literal cleared trace equations are
\[
 w_0^2-2w_0\tau-2w_1^2=0,\qquad (1+q)\mathfrak b=\gamma^2\tau^2.
\]
By Lemma~\ref{lem:residues}, $W\ne0$. The first equation therefore
gives $w_0\ne0$ and $\tau=1/\Tr(W^{-1})\ne0$; indeed $w_0=0$
would give $w_1=0$. Hence
$\mathfrak b=\gamma^2R/\Tr(W^{-1})^2$. The factor $\gamma$ is a unit at every consumed contact prime. If
$\ell=\max_{v\mid p}v(W)\le0$, the trace inequality gives
$v_p(\mathfrak b)\le\nu+2\ell$.
This valuation is odd; the retained pole obstruction forces it positive.
Consequently $\nu>0$ and $\ell>-\nu/2$.
At a place attaining $\ell$, the factor valuations are nonpositive
and add to $\ell$. Each factor valuation is therefore at least $\ell$,
and each of its term valuations is strictly greater than $-\nu/2$.

The resulting lower bounds, before integer rounding, are
\[
\begin{array}{c|c|c}
\text{test}&\text{strict lower bound}&\text{integer order}\\\hline
C_2,C_3,C_4,G&9\nu/4&3\\
I&9\nu/8&2\\
C_1&3\nu/8&1\\
D_H&3\nu/4&2\\
A_2&5\nu/4&1\\
A_1,A_3,A_4&3\nu/4&1\\
B_0,B_1,B_2,B_3,B_4&7\nu/8&1
\end{array}
\]
Here $\nu$ is a positive odd integer. The order two for $D_H$ uses
its even valuation. In fact each of the nine reciprocal-abscissa
tests has order at least two; order one suffices for its subsequent
formal-group use. Zero tests satisfy all the stated inequalities.
\end{proof}

\section{Integrality before the formal congruences}
At a contact place selected by Proposition~\ref{prop:extract}, the
points $H,U_i$ are formal in the primary completion, and $H-P,V_i$
are formal in the conjugate completion. The reciprocal-abscissa
tests have positive even valuation, so every nonidentity point has
formal depth at least one. In the parameter $z=-x/y$, write
\[
 \log_E(Q)=z(Q)\mathcal A(Q),\qquad
 \mathcal A(Q)\equiv1\pmod{p^{4v(z(Q))}}.
\]
The value of the unit is taken to be one at the identity. This is
the constant-term formal logarithm estimate proved in~\cite{Logs}.

The integral group identities
\[
                  [m]U_i=[2u_i]H,\qquad
                  [m-2]V_i=[2(u_i-n_i)](H-P)
\]
kill the torsion in $H=[m]R_0+T_0$. Their full logarithmic ratios
are $u_i/h$ and $(u_i-n_i)/(h-1)$, where
$(n_1,n_2,n_3,n_4)=(n,a,t,b)$. At least one of $h,h-1$ is a
$p$-adic unit, so at least one full ratio is integral. The formal
units preserve the valuation of its plain parameter ratio.
The positive-order test $C_i$ transfers integrality to the other
plain ratio, and then to the other full ratio.

Only after this transfer may the relative errors of order four
be used as absolute congruences. They give
\[
 L_i\equiv u_i/h\pmod{p^4},\qquad
 L_i'\equiv(u_i-n_i)/(h-1)\pmod{p^4}.
\]
Since $n_i=h(u_i/h)-(h-1)((u_i-n_i)/(h-1))$, the high cross tests
give $L_i\equiv n_i\pmod{p^3}$ for $i>1$. The first cross test
gives $L_1\equiv n\pmod p$, while $I$ gives $L_1\equiv\beta\pmod p$.
These arguments include identity numerator points, half-integral $h$,
and inert primes, where conjugation preserves the unique valuation.

\section{Global soundness for sufficiently large parameters}\label{sec:soundness}
The scalar bounds~\eqref{eq:scalar} hold independently of the
comparison logarithms. In particular $c\to\infty$ uniformly with $N$,
and $B\ge c-1$.

Let
$\Delta=t^2-a(a-2)tb+(a-2)^2b^2-1$.
The local comparison shows that each third power of a comparison
radical divides $\Delta$. If $\Delta\ne0$, then
$|\Delta|\le14B^4$, and \eqref{eq:heighttransfer} gives
\[
                  B^2\le C14^{7/18}B^{14/9}+M_2.
\]
Since $14/9<2$, this bounds $B$ by a fixed constant. Equation
\eqref{eq:scalar} gives $B\ge c-1\to\infty$ uniformly as
$N\to\infty$, a contradiction. Thus $\Delta=0$ eventually.
A stronger form of the integral Lucas descent gives
\begin{equation}\label{eq:stronglucas}
              \log B\ge(c-5)\log(c-2)\qquad(c\ge8).
\end{equation}
Here is the complete elementary argument. Put $A_0=|a|$ and
$w=\operatorname{sgn}(a)(a-2)b$. Then
$t^2-A_0tw+w^2=1$. For
$\eta=(A_0+\sqrt{A_0^2-4})/2$, every positive norm-one unit of
$\Z[\eta]$ is a power of $\eta$. Indeed multiply a unit by a
power of $\eta$ to place it in $[1,\eta)$; its coefficient of
$\eta$ is $(v-v^{-1})/(\eta-\eta^{-1})$, an integer in $[0,1)$,
so it is zero and $v=1$.
Let $U_0=0$, $U_1=1$, and $U_{j+1}=A_0U_j-U_{j-1}$ for $j\ge1$.
The unit identities show that $|w|=U_j$ and
$|t|\in\{U_{j-1},U_{j+1}\}$ for some $j\ge1$.
Modulo $A_0-2$ one has $U_j\equiv j$; modulo $A_0+2$ one has
$U_j\equiv(-1)^{j-1}j$. Since $w$ is a nonzero multiple of
$|a-2|$, it follows that $j\ge A_0-2$ if $a>0$ and
$j\ge A_0+2$ if $a<0$. Induction gives
$U_j\ge(A_0-1)^{j-1}$ for $j\ge1$. Therefore
$|t|\ge(A_0-1)^{A_0-4}$. Finally $A_0\ge c-1$, which proves
\eqref{eq:stronglucas}.

If $n-\beta\ne0$, let $J$ be its nonzero reduced numerator.
The local input comparisons make both radicals divide $J$,
while $|J|\le2B\exp(h(\beta))$. For $B^2\ge2M_2$,
\eqref{eq:heighttransfer} gives
\[
 \log B\le\frac65\log(2C)+\frac75\log2+\frac75h(\beta)
 \le C'+\frac75H\ell^{-7L_N/9}c^{7/9}.
\]
The last expression is $o(c)$ as $c\to\infty$, uniformly for
$N\ge1$. The Lucas bound $(c-5)\log(c-2)$ grows faster than $c$.
Since $c\ge\ell^{L_N}\to\infty$ uniformly with $N$, these two
inequalities are incompatible for all sufficiently large $N$.
Hence $n=\beta\in\Z$. If $X=u/v$ is reduced with $v>0$, then
$v^2\mid51$, and squarefreeness of $51$ gives $v=1$.
All bounds depend only on fixed constants. Since $L_N$ is increasing,
a single $N_0$ works for every larger $N$. Its dependence on the
aggregate additive constant is ineffective.

\section{Simultaneous completeness}
Fix an integer input and $N\ge1$. Choose a real neighborhood of a
nonzero reference point $[h_*]P$, with $64\mid h_*$, on which
both values of $|z(H)|$ are bounded below. Then choose a sufficiently
small real neighborhood of the identity for $C_S$. It gives
$|Z_0|,|Z_1|\le1$. For nonzero $H\in\Z P$ the positive-square
abscissa bound also gives $|\Nm(f(H))|\le1/7$.
For $R_s\ge1$, equations \eqref{eq:scalarnorm} imply
\[
 A\le2R_s^{-3},\qquad
 W_s\le4R_s^{-10}+\frac1{49}R_s^{-3}<5R_s^{-3}.
\]
Therefore
\[
 0<b_s\le25\frac{R_s+9792}{R_s^5}<17\qquad(R_s\ge10).
\]
The displayed rational function is decreasing on $(0,\infty)$, and
its value at ten is less than $17$. Since $r_s\ge1$,
\[
 R_s\ge\frac{149^{2N+1}}{1+149^2}
       \ge\frac{149^3}{1+149^2}>10.
\]
Thus the required real preparation holds for every $N\ge1$ with
explicit margins.

The finite preparation uses no new arithmetic theorem. For integer
input, $r_s$ is a positive integer and
\[
 R_s(R_s+9792)=
 \frac{A_Nr_s\bigl(A_Nr_s+9792(1+\ell^2r_s)\bigr)}
                    {(1+\ell^2r_s)^2}.
\]
The numerator is coprime to the displayed denominator, so its reduced
denominator is square. At the three preparation primes the elementary
input congruences give
\[
 v_2(R_s)\in\{2,3\},\qquad v_3(R_s)=1,
                       \qquad v_{17}(R_s)=0.
\]
In particular $r_s\equiv15\pmod{17}$ and
$1+149^2r_s\equiv3\pmod{17}$, so the normalizer denominator is
also a unit at $17$. Thus
$R_s(R_s+9792)=R_s^2(1+9792/R_s)$ is a local square at all three primes.
Choose the shared index $h_0$ sufficiently divisible at the finite
support of this rational number, and then choose the index
$c=a-1$ still more divisible. The points $H=[h_0]P$ and $C_S=[c]P$
are then formal at the required finite places, and all components of
$f(H)$ and $Z$ have the needed arbitrarily large valuations. Make
$W_s$ integral at every prime where the displayed numerator has odd
valuation, and deep enough at $2,3,17$ to supply the pole bounds.
At every other prime the valuation of $b_s$ is even. These conditions
give a square reduced denominator and all pole preparation hypotheses.
The source's progression and compact-real-group argument supplies the
chosen real neighborhoods at the same time, together with the retained
nonzero chart guards. It also permits $c$ to be arbitrarily large and
to satisfy the retained divisibility conditions for the Lucas construction.

After $S$ is fixed, continuity preserves this finite scalar preparation
and its guards for all $h$ in one fixed congruence class modulo a
sufficiently divisible integer $L$. Retain the source's integral Lucas
construction, choose the four root labels, and fix both positive scales
and compact neighborhoods for all comparison tests before selecting
contact primes. The additional contact primes avoid the finite scalar
preparation support. Hence these later congruences do not conflict with
any scalar congruence.

Choose a fixed positive integer $D_0$ clearing all contact coefficients
$D_0^2\theta/a_*$ and sufficiently divisible at the prepared primes,
including $\Sigma$, and the prime divisors of the scales. In the prime
pattern write
\[
                         \lambda=\frac{u}{D_0^2w^2}.
\]
The zero contact makes $u$ one of the selected primes, outside that
finite set. Thus $v_p(\lambda)\le-2v_p(D_0)$ there. Make this upper
bound low enough that $t_*=a_*\lambda$ has valuation below zero and
all nonzero contacts. Then $v_p(\delta)=v_p(t_*)$, and
\begin{equation}\label{eq:preparedsquare}
 R=\prod_{\theta\in\Theta_2}
                 \left(1+\frac{1-\theta}{\delta}\right)
       \longrightarrow1
\end{equation}
in every prepared completion. The weight $q=R^{-1}-1$ tends to zero.
Increasing the fixed denominator divisibility makes its valuation
uniformly arbitrarily large. Thus all weighted test terms are small
at every prepared place, and $\Tr(W^{-1})$ is arbitrarily close to two.

The factors in~\eqref{eq:preparedsquare} are arbitrarily close to one.
Requiring congruence to one modulo eight at $2$, and modulo $p$ at an odd
prepared prime $p$, makes their product a local square. The pole input
is $\gamma^2R/\Tr(W^{-1})^2$, so it is also a square at $2,3,17$.
The fixed factor has valuations
\[
 v_2(\gamma)=3,\qquad v_3(\gamma)=v_{17}(\gamma)=1.
\]
Since the reciprocal trace is close to $1/2$, the pole input has
valuations $4,2,2$ at those three primes. These satisfy all retained
preparation inequalities.

This preparation is valid for any fixed nonzero rational scales.
The prime-pattern input in \cite{Selection,Dependencies} applies to
arbitrary fixed distinct integer coefficients. After clearing
$D_0^2\theta/a_*$ and removing repeated coefficients, it applies
here without a square restriction on $a_*$. All denominator
conditions are fixed before prime selection.

Select the contact primes with the inherited opposite-order separation.
Take $H=[h]P$, $U_i=[hn_i]P$ and $V_i=[(h-1)n_i]P$.
The common-index congruences preserve the scalar progression and make
$H,U_i$ formal at one place and $H-P,V_i$ formal at its conjugate.
Only natural depth one is needed on either branch; the source's stronger
primary preparation is also permitted. Both full logarithmic ratios
are exactly $n_i$. Plain-parameter errors therefore have order at least
four. All $C_i$, and $I$ since $n_1=\beta$, have valuation at least four,
and the Lucas equation for the chosen labels gives $v(G)\ge4$.
Every reciprocal-abscissa test has valuation at least two.
At the selected place $v_p(q)=-1$, so every nonconstant term of every
packed factor has positive valuation. Thus $W\equiv1$ there. At the
opposite place, each factor has valuation at most zero, with residue one
when its valuation is zero. The trace of $W^{-1}$ is a unit with residue
one or two. The comparison pole input has valuation one at every
selected contact prime.

The exact prime-pattern identity is
\[
 R=\frac{\displaystyle a_*^{10}\prod_{\theta\in\Theta_2}
          (u-(\theta/a_*)D_0^2w^2)}
        {(D_0^2w^2)^{10}\delta^{10}}.
\]
Every factor outside the selected prime numerators is a rational
square, including $a_*^{10}$ even when $a_*$ is not itself a square.
Thus at every prime outside those numerators the valuation is even. Multiplication by $\gamma^2$ and division by the squared trace preserve this parity. This argument includes primes dividing $\delta$;
no unit assertion about $\delta$ is made away from contacts.
Equation~\eqref{eq:preparedsquare} and the preparatory divisibility
give the required local squareclasses and valuation bounds at $2,3,17$.

For the real preparation, write $r=\delta^2$ and let
$\alpha_1,\ldots,\alpha_5>0$ be the five squared distances of the
contact pairs from one. Then
\[
 R=\prod_{j=1}^{5}(1-\alpha_j/r),\qquad
 q\sim\frac{\alpha_1+\cdots+\alpha_5}r>0
 \quad\text{as }r\longrightarrow+\infty.
\]
On the fixed compact real neighborhoods the tests remain bounded, so
$W\to1$ and $\Tr(W^{-1})\to2$. Consequently
\[
              \mathfrak b\longrightarrow\frac{\gamma^2}{4}
                              =\frac{41616}{167281}\in(0,17).
\]
All factor values are positive at both real embeddings for sufficiently
large $\lambda$, because $q>0$ and all test powers are even.
The simultaneous selection supplies the required real size after all
fixed data have been chosen. The pole-realization input supplies its
witnesses, and the exact power identities fill the quadratic circuit.
This proves completeness for every integer input and every $N\ge1$.

\section{Exact quadratic specification}
Retain three scalar inverse weights $s_1,s_2,s_4$ for each detector,
with the quadratic equations
\begin{equation}\label{eq:inverseweights}
                 qs_1=1,\qquad s_2=s_1^2,\qquad s_4=s_2^2.
\end{equation}
Thus $q\ne0$ and $s_j=q^{-j}$ for $j=1,2,4$. The actual factor
coordinates $T_j$ are specified by the following quadratic equations:
\begin{equation}\label{eq:clearedfactors}
\begin{split}
 s_4T_1&=s_4+qC_2^2+s_2D_H^2+(B_1^2)^2,\\
 s_4T_2&=s_4+qG^2+s_2A_3^2+(B_0^2)^2,\\
 s_4T_3&=s_4+qI^4+s_1A_2^2+s_2A_4^2+(B_3^2)^2,\\
 s_4T_4&=s_4+qC_3^2+s_2C_1^4+(B_4^2)^2,\\
 s_4T_5&=s_4+qC_4^2+s_2A_1^2+(B_2^2)^2.
\end{split}
\end{equation}
Every test square on the right is a stored field value. Squaring
the stored $B_i^2$ is quadratic in its two rational components and
needs no output for $B_i^4$. Multiplication by $q^4$ proves exact
equivalence with the five factor definitions above. The products of
the factors and the pole normalization are unchanged by this clearing.

The three contact outputs $r,r_2,v$ are specified by
\eqref{eq:contactcircuit}, using the already retained inverse weight
$s_1$. Their equations involve only products of two stored quantities.
No coordinate for $P_5(r)$ is required. Lemma~\ref{lem:contactcircuit}
proves the exact rational interpretation and treats the $r=0$ stratum
without division by $r$.
The scalar normalizer is the quadratic equation
$R_s+22201r_sR_s-A_Nr_s=0$.
The scalar chain is $R_2=R_s^2$, $R_4=R_2^2$.
The two scalar norm rows use $R_4$, as in~\eqref{eq:scalarnorm}.

Only $C_1^4$ and $I^4$ are retained as separate field fourth powers.
The five $B_i^4$ values occur directly as squares of stored field
coordinates in \eqref{eq:clearedfactors}. All remaining point and
scalar equations are fixed by the literal coefficient tables.
The exact count is
\[
\begin{array}{lr}\toprule
\text{point, scalar, contact, ratio, and weight coordinates}&190\\
\text{two shared fourth powers}&4\\
\text{five factors and intermediate products for two detectors}&32\\\midrule
\text{total coordinates, including the input}&226\\\midrule
\text{point, scalar, contact, ratio, and weight quadratic rows}&160\\
\text{shared fourth-power rows}&4\\
\text{factor and product rows}&36\\\midrule
\text{total quadratic rows}&200\\\bottomrule
\end{array}
\]

The complete tables fix every coordinate, equation, coefficient, and
normalization; the input is coordinate one. For the common point,
the six primitive rational component equations are replaced by the
fixed rational row combinations in the certificate, whose determinant
is $-1/362963666515968$. The producer clears denominators and makes
each resulting row primitive before squaring. The row change is
nonsingular and preserves the common zero set.
The independent coefficient checker imports no producer. It computes
the elementary symmetric sums of the five squared contact distances,
checks the five roots of $P_5$, and verifies the general polynomial
identity in Lemma~\ref{lem:contactcircuit}. It also verifies the
original normalization algebraically after multiplication by $q$;
the discrepancy is exactly a multiple of $qs_1-1$. The nonzero
constant term of $P_5$ certifies the exceptional-stratum argument.
The checker inspects all occurrences of every omitted polynomial
coordinate in the pinned source. It reconstructs the 4 changed
contact rows and compares the other 196 rows exactly with that
source. It independently expands all 200 quadratic squares and
checks every coefficient, every coordinate, the degree, the content,
and the coefficient-height formula.

The reconstruction dependencies retain separate exact certificates
for the other algebraic components. The common-point certificate
verifies twelve polynomial identities, including elimination without
division and translation on every exceptional stratum. The factor
certificates check the residue classes for all odd contact orders,
the resulting comparison precisions, and the completeness depths.
The contact-set certificate checks the exceptional primes and the
height-transformation matrices. The scalar certificate checks the
forced prime, its nonsquare residue tests, and the extraction bounds.
These finite certificates support the algebraic specification; the
infinite arithmetic arguments and their premises are the proofs above.

The expanded quartic has 2836 terms before powers of $A_N$ belonging
to the same coordinate monomial are collected, and 2835 monomials
after collection. Its unique coefficient involving $A_N^2$ is
$A_N^2+1$, at the square of the actual scalar contact-product coordinate.
Every other coefficient is a nonzero constant or a nonzero constant
multiple of $A_N$. Consequently no support cancellation occurs for
positive $A_N$. The largest constant coefficient $C_0$ has 1520 bits,
the largest linear multiplier is $44402$, and exact comparison
gives
\[
 C_0>44402^2,\qquad
 149^{210}+1<C_0<149^{214}+1.
\]
Thus $H(P_N)=\max(C_0,A_N^2+1)$ for every $N\ge1$ and the asserted
exact height holds for $N\ge53$. This is a coefficient statement;
it does not make the defining threshold effective. The producer
\nolinkurl{build_inverse_contact_v21.py} and the independent checker
\nolinkurl{check_inverse_contact_v21.py} reconstruct all 200 rows and
every coefficient. The full coefficient tables and certificate are in
\nolinkurl{computation/quartic/family226_v21}.

\begin{thebibliography}{99}\small\setlength{\itemsep}{2pt}
\bibitem{Dependencies} Project repository,
\emph{Arithmetic dependencies of the defining-polynomial constructions},
\nolinkurl{arithmetic_dependencies_20261108_v8.tex}.
\bibitem{Source} Project repository,
\emph{A family of quartics in 297 variables},
\nolinkurl{quartic297_20261109_v12.tex}, especially Sections 2--6.
\bibitem{Weights} Project repository,
\emph{Common local weights and quartics in 297 and 342 variables},
\nolinkurl{shared_local_weights_20261109_v12.tex}.
\bibitem{Even} Project repository,
\emph{An explicit quartic in 322 variables},
\nolinkurl{quartic322_20261109_v13.tex}.
\bibitem{Contacts} Project repository,
\emph{Height Bounds from Disjoint Sets of Rational Contacts},
\nolinkurl{contact_height_amplification_20261109_v13.tex}.
\bibitem{Aggregate} Project repository,
\emph{The height of an aggregate elliptic trace},
\nolinkurl{aggregate_trace_height_20261108_v9.tex}.
\bibitem{NewHeight} Project repository,
\emph{Evaluation of Differentials and Heights of Modular Jacobians},
\nolinkurl{evaluation_of_differentials_20261109_v17.tex},
Definition~6.1 and Theorem~6.2.
\bibitem{Multiplicity} Project repository,
\emph{Multiplicity of modular factors and quadratic contact-height bounds},
\nolinkurl{multiplicity_and_contact_heights_20261110_v19.tex},
Theorem~4.1.
\bibitem{Logs} Project repository,
\emph{Uniform Rational Approximations to an Elliptic Logarithm},
\nolinkurl{elliptic_logarithms_20261108_v8.tex}.
\bibitem{Selection} Project repository,
\emph{Prime selection with fixed congruences and separated reduction orders},
\nolinkurl{prime_selection_20261109_v12.tex}.
\end{thebibliography}
\end{document}
